\hypertarget{exact-sequence-polyphase-and-tap-laboratory}{%
\section{Exact sequence, polyphase and tap
laboratory}\label{exact-sequence-polyphase-and-tap-laboratory}}



This chapter connects the nonlinear family to exact sequences and finite
filter algebra. It supplies reusable code without replacing the analytic
proofs in \path{routes/asymmetric_dilation.md}. All mathematical
sequence indices begin at zero. All finite arrays are zero extended
beyond their displayed support. The catalog contains explicit candidates
for OEIS lookup. This chapter does not assert an OEIS accession or
historical priority for an unsearched row.

\hypertarget{the-integer-polynomial-object}{%
\subsection{The integer polynomial
object}\label{the-integer-polynomial-object}}

Write the unique odd formal solution as

\[
T_{q,r}(z)=\sum_{n\ge0}c_n(q,r)z^{2n+1}
          =\sum_{n\ge0}A_n(q,r)\frac{z^{2n+1}}{(2n+1)!}.
\]

The two coefficient systems are related by \(A_n=(2n+1)!c_n\). In
particular, \(A_n\) is the derivative of order \(2n+1\) at zero. The
coefficient \(c_n\), the derivative \(A_n\), and the full EGF sequence
\(a_m\) with \(a_{2n+1}=A_n, a_{2n}=0\) are different objects. None of
the code silently substitutes one for another.

The nonlinear differential equation gives the recurrence

\[
\boxed{
A_0=1,\qquad
A_n=\sum_{i=0}^{n-1}\binom{2n}{2i+1}
 q^{2i+1}r^{2(n-1-i)+1}A_i A_{n-1-i},\quad n\ge1.}
\]

\textbf{Proof.} Multiply the ordinary Taylor recurrence by \((2n+1)!\).
The factor attached to \(A_iA_j\), where \(i+j=n-1\), is
\((2n)!/[(2i+1)!(2j+1)!]=\binom{2n}{2i+1}\). This also proves the
recurrence directly in the EGF convention. \(\square\)

\textbf{Polynomial theorem.} For every \(n\), \(A_n\in\mathbb Z[q,r]\),
all its nonzero coefficients are positive integers, and
\(A_n(q,r)=A_n(r,q)\). Every monomial \(q^u r^v\) satisfies

\[
u,v\ge n,\qquad u\equiv v\equiv n\pmod2,\qquad u+v\le n(n+1).
\]

\textbf{Proof.} The recurrence is an induction using positive integer
binomial coefficients. Swapping \(q,r\) and replacing \(i\) by \(n-1-i\)
leaves the sum unchanged. If child indices are \(i,j\), their exponent
pairs receive the shifts \((2i+1,2j+1)\). The lower bounds follow from
\(i+j+1=n\). The two parities are the parity of \(i+j+1\), also \(n\).
For the degree bound, a recurrence term has total degree at most

\[
2n+i(i+1)+j(j+1)\le n(n+1),\qquad i+j=n-1,
\]

since \(i(i+1)+j(j+1)\) is largest at \((i,j)=(0,n-1)\) or its swap. The
base \(A_0=1\) satisfies every assertion. \(\square\)

The first four polynomials are

\[
\begin{aligned}
A_0={}&1,\\
A_1={}&2qr,\\
A_2={}&8q^2r^2(q^2+r^2),\\
A_3={}&q^3r^3\bigl[48(q^2+r^2)(q^4+r^4)+80q^2r^2\bigr].
\end{aligned}
\]

Consequently each integer pair \((q,r)\) gives an integer sequence for
all \(n\), independently of whether its formal Taylor series converges.
For example, the first values of \(A_n(1,2)\) are
\(1,4,160,35200,38041600\).

\hypertarget{rational-feedback-banks-also-admit-integer-sequences}{%
\subsubsection{Rational feedback banks also admit integer
sequences}\label{rational-feedback-banks-also-admit-integer-sequences}}

For rational channel tuples \((h_t,q_t,r_t)\), use

\[
F'(z)=1+\sum_t h_tF(q_tz)F(r_tz),\qquad F(0)=0.
\]

The preceding recurrence is summed over channels with the factor
\(h_t\). In the universal polynomial over the weights and contractions,
every term of \(A_n\) has weight degree exactly \(n\) and total
contraction degree at most \(n(n+1)\). The weight statement follows from
one new weight at each recurrence step and \(1+i+j=n\); the contraction
bound uses the same induction.

Let \(H\) be a common denominator for all channel weights, and \(D\) a
common denominator for all channel contractions. Then

\[
\boxed{B_n=H^nD^{n(n+1)}A_n\in\mathbb Z\quad\hbox{for all }n.}
\]

This is a proved normalization, even with signed rational weights. It is
usually not the smallest possible denominator clearing. The catalog
records both \(A_n\) and \(B_n\), and the exact multiplier, so that a
proposed sequence lookup can state its normalization honestly.

\hypertarget{a-sharp-distinction-between-sequences-and-analytic-germs}{%
\subsection{A sharp distinction between sequences and analytic
germs}\label{a-sharp-distinction-between-sequences-and-analytic-germs}}

The entire-function theorem uses \(0<q,r<1\). Specializations with large
positive integers produce valid formal and integer sequences but do not
necessarily produce analytic functions near zero.

\textbf{Zero-radius theorem.} If \(q>0\) and \(r>1\), the formal series
for \(T_{q,r}\) has radius of convergence zero. The same holds after
swapping the parameters.

\textbf{Proof.} Every coefficient is positive. In the ordinary
recurrence keep only the term \(i=0\). This gives

\[
c_n\ge\frac{q r^{2n-1}}{2n+1}c_{n-1}
\quad\Longrightarrow\quad
c_n\ge\frac{q^n r^{n^2}}{3\cdot5\cdots(2n+1)}.
\]

Since the denominator is at most \((2n+1)^n\),

\[
c_n^{1/(2n+1)}\ge
q^{n/(2n+1)}r^{n^2/(2n+1)}(2n+1)^{-n/(2n+1)}\longrightarrow\infty.
\]

The Cauchy--Hadamard criterion therefore gives radius zero. \(\square\)

More generally, a positive feedback bank with a channel satisfying
\(h>0,q>0,r>1\) has the same obstruction, replacing \(q^n\) by
\((hq)^n\). Positivity prevents cancellation. No conclusion of this form
is asserted for a bank with signed weights.

\begin{longtable}[]{@{}
  >{\raggedright\arraybackslash}p{(\columnwidth - 4\tabcolsep) * \real{0.3333}}
  >{\raggedright\arraybackslash}p{(\columnwidth - 4\tabcolsep) * \real{0.3333}}
  >{\raggedright\arraybackslash}p{(\columnwidth - 4\tabcolsep) * \real{0.3333}}@{}}
\toprule\noalign{}
\begin{minipage}[b]{\linewidth}\raggedright
Parameters or bank
\end{minipage} & \begin{minipage}[b]{\linewidth}\raggedright
First five \(A_n\)
\end{minipage} & \begin{minipage}[b]{\linewidth}\raggedright
Mathematical status
\end{minipage} \\
\midrule\noalign{}
\endhead
\bottomrule\noalign{}
\endlastfoot
\((1,1)\) & \(1\), \(2\), \(16\), \(272\), \(7936\) & Known tangent derivatives \\
\((0,1)\) & \(1\), \(0\), \(0\), \(0\), \(0\) & Known linear boundary \\
\((1,-1)\) & \(1\), \(-2\), \(16\), \(-272\), \(7936\) & Known hyperbolic tangent
derivatives \\
\((1,2)\) & \(1\), \(4\), \(160\), \(35200\), \(38041600\) & Integer sequence, radius
zero \\
\((1/2,1/2)\) & \(1\), \(1/2\), \(1/4\), \(1/8\), \(1/16\) & Known scaled hyperbolic sine \\
\((1/4,3/4)\) & \(1\), \(3/8\), \(45/256\), \(21465/262144\), \(39833775/1073741824\) &
Entire asymmetric family, order one \\
\end{longtable}

The complete catalog contains \textbf{18 rows}, each through \(n=12\):
three classical or degenerate integer anchors, eight positive integer
formal specializations, four rational contractions, one entire edge case
of infinite order, and two normalized feedback banks. Integer
integrality and rational clearing are established by the preceding
proofs, not merely inferred from the finite sample.

\hypertarget{the-uncontracted-boundary-is-entire-of-infinite-order}{%
\subsubsection{The uncontracted boundary is entire of infinite
order}\label{the-uncontracted-boundary-is-entire-of-infinite-order}}

\textbf{Boundary theorem.} For \(0<q<1\), the normalized formal solution
of

\[
T'(z)=1+T(qz)T(z),\qquad T(0)=0
\]

extends to an entire function of infinite order. The swapped boundary
has the same property.

\textbf{Proof of entire extension.} The tangent-coefficient majorant
gives an initial holomorphic germ. Suppose \(T\) is known on \(|z|<R\).
The function \(T(qz)\) is known on \(|z|<R/q\). On this larger disk
solve the linear ODE \(U'=1+T(qz)U\), \(U(0)=0\). Explicitly,

\[
U(z)=e^{B(z)}\int_0^z e^{-B(w)}\,dw,
\qquad B(z)=\int_0^zT(qw)\,dw.
\]

The integrals are path independent on the disk. Linear-ODE uniqueness
shows \(U=T\) on the previous disk, and there \(T(qz)\) equals the newly
extended function at \(qz\). Thus the same nonlinear equation holds on
the enlarged disk. Iterating \(R\mapsto R/q\) exhausts the plane.

\textbf{Proof of infinite order.} Positive coefficients imply
\(T(x)\ge x\) for \(x>0\) and \(M(x)=T(x)\). The logarithmic derivative
obeys

\[
\frac{d}{dx}\log T(x)=\frac1{T(x)}+T(qx)\ge qx.
\]

Integrating from one, with \(T(1)\ge1\), gives
\(T(x)\ge\exp(q(x^2-1)/2)\) for \(x\ge1\). For \(R\ge2/q\), integrate
the same logarithmic derivative over \([R/2,R]\). Since \(T(R/2)\ge1\),

\[
\boxed{
\log M(R)\ge\frac R2\exp\left(\frac{q^3R^2}{8}-\frac q2\right).}
\]

Therefore \(\log\log M(R)/\log R\to\infty\), which proves infinite
order. \(\square\)

For positive parameters this gives a complete analytic/formal phase
chart:

\begin{longtable}[]{@{}
  >{\raggedright\arraybackslash}p{(\columnwidth - 2\tabcolsep) * \real{0.5000}}
  >{\raggedright\arraybackslash}p{(\columnwidth - 2\tabcolsep) * \real{0.5000}}@{}}
\toprule\noalign{}
\begin{minipage}[b]{\linewidth}\raggedright
Parameter region
\end{minipage} & \begin{minipage}[b]{\linewidth}\raggedright
Normalized solution
\end{minipage} \\
\midrule\noalign{}
\endhead
\bottomrule\noalign{}
\endlastfoot
\(0<q,r<1\) & Entire, finite order defined by \(q^\rho+r^\rho=1\) \\
One parameter equals one, the other lies in \((0,1)\) & Entire, infinite
order \\
\(q=r=1\) & Tangent, meromorphic with its usual poles \\
\(q>0,r>0,\max(q,r)>1\) & Formal solution with radius zero \\
One parameter equals zero & Linear solution \(z\), regardless of the
other scale \\
\end{longtable}

The finite-order formula applies in the open contraction square. It does
not assign a finite order to the edge, where its defining scale equation
has no positive root.

\hypertarget{exact-root-filter-sections}{%
\subsection{Exact root-filter
sections}\label{exact-root-filter-sections}}

For an EGF \(G(z)=\sum_{m\ge0}a_m z^m/m!\), an integer base \(b\ge1\), a
residue \(0\le j<b\), and \(\omega=e^{2\pi i/b}\), define

\[
\boxed{
\mathcal P_{b,j}G(z)=\frac1b\sum_{k=0}^{b-1}\omega^{-jk}G(\omega^kz)
=\sum_{m\equiv j\,(b)}a_m\frac{z^m}{m!}.}
\]

This holds as a formal series and wherever the original series
converges. Indeed the coefficient multiplier is the geometric sum
\(b^{-1}\sum_{k=0}^{b-1}\omega^{k(m-j)}\), equal to one or zero
according to the residue. It proves all the projector identities

\[
\sum_{j=0}^{b-1}\mathcal P_{b,j}=I,\qquad
\mathcal P_{b,j}\mathcal P_{b,l}=\delta_{jl}\mathcal P_{b,j}.
\]

The full-position section and the decimated lane are distinct arrays. A
section keeps the original positions, inserting zeros outside the
residue. A lane \(x_j[n]=a_{bn+j}\) changes the index to \(n\).

For an \textbf{ordinary} generating function \(X(z)=\sum_ma_m z^m\),
define \(X_j(w)=\sum_na_{bn+j}w^n\). Then

\[
X(z)=\sum_{j=0}^{b-1}z^jX_j(z^b).
\]

For an EGF the corresponding exact reconstruction is

\[
G(z)=\sum_{j=0}^{b-1}z^j
 \sum_{n\ge0}\frac{a_{bn+j}}{(bn+j)!}(z^b)^n.
\]

One must retain \((bn+j)!\). Replacing it by \(n!\) gives a different
function. The root-filter section identity has no such factorial
ambiguity.

\hypertarget{ordinary-taps-egf-multiplication-and-differentiation}{%
\subsection{Ordinary taps, EGF multiplication and
differentiation}\label{ordinary-taps-egf-multiplication-and-differentiation}}

Let \(h_0,\ldots,h_{L-1}\) be a finite tap vector and let \(a_m=0\) for
\(m<0\). Ordinary taps form

\[
b_m=\sum_{t=0}^{L-1}h_ta_{m-t}.
\]

For OGFs this is multiplication by \(H(z)=\sum_th_tz^t\):
\(B(z)=H(z)X(z)\). For EGFs it is a sum of repeated zero-constant
antiderivatives:

\[
\boxed{B_{\rm EGF}(z)=\sum_{t=0}^{L-1}h_t I^tG(z),\qquad
IG(z)=\int_0^zG(w)\,dw.}
\]

\textbf{Proof.} The EGF of the delayed sequence is
\(\sum_{m\ge t}a_{m-t}z^m/m!=I^tG(z)\). Summing taps gives the formula.
Conversely, \(G^{(t)}\) has coefficient sequence \(a_{m+t}\), so
derivatives sample future coefficients and do not implement these past
delays. \(\square\)

For two EGFs \(G_a,G_h\), ordinary function multiplication instead gives
the \textbf{Hurwitz product}

\[
G_a(z)G_h(z)=\sum_m\left[
\sum_{t=0}^m\binom mt h_ta_{m-t}\right]\frac{z^m}{m!}.
\]

This is a binomial convolution. For example, with \(a_m=1\) and taps
\(h=(1,1)\), ordinary taps give \(b_2=2\), while the Hurwitz product
gives \(b_2=3\). This explicit negative control catches a common
normalization error.

\textbf{Intrinsic nonlinear feedback taps} are the parameters
\((h_t,q_t,r_t)\) inside \(F'=1+\sum_th_tF(q_tz)F(r_tz)\). They change
the function's defining equation and its derivative recurrence. Ordinary
FIR taps are a linear transformation of an already defined coefficient
sequence. Both are useful, but they operate on different objects. The
intrinsic-bank growth theorem belongs to
\path{routes/asymmetric_dilation.md}.

\hypertarget{the-exact-finite-polyphase-tap-matrix}{%
\subsection{The exact finite polyphase tap
matrix}\label{the-exact-finite-polyphase-tap-matrix}}

For an OGF sequence \(a_m\), define lanes \(x_l[n]=a_{bn+l}\) and
\(X_l(w)=\sum_nx_l[n]w^n\). Apply finite ordinary taps and obtain
\(y_j[n]=b_{bn+j}\). Then the exact matrix identity is

\[
\boxed{Y_j(w)=\sum_{l=0}^{b-1}H_{jl}(w)X_l(w),\qquad
H_{jl}(w)=\sum_{\substack{0\le t<L\\t+l\equiv j\,(b)}}
h_t w^{(t+l-j)/b}.}
\]

Every exponent is a nonnegative integer. It records the carry across
phase boundaries, rather than erasing it.

\textbf{Proof.} A term \(h_t a_{bn+j-t}\) is in phase \(l\) exactly when
\(t+l\equiv j\pmod b\). Write \(d=(t+l-j)/b\). That term is
\(h_t x_l[n-d]\), whose OGF is \(h_tw^dX_l(w)\) by zero extension.
Summing over taps and phases gives the matrix. \(\square\)

For two phases and two taps,

\[
\begin{pmatrix}Y_0\\Y_1\end{pmatrix}
=\begin{pmatrix}h_0&h_1w\\h_1&h_0\end{pmatrix}
\begin{pmatrix}X_0\\X_1\end{pmatrix}.
\]

The upper-right factor \(w\) is necessary: a delay from odd phase to
even phase crosses one block. A diagonal-only phase operation cannot
implement this filter. The code constructs the sparse matrix for any
finite base and tap vector and checks that recombination gives the full
direct convolution.

\hypertarget{reuse-and-evidence}{%
\subsection{Reuse and evidence}\label{reuse-and-evidence}}

From the project directory, standard Python is enough:

\begin{verbatim}
python family_lab.py coefficients --q 1/4 --r 3/4 --n 8
python family_lab.py catalog --n 12 --output-dir examples
python checks/family_lab_checks.py
\end{verbatim}

The reusable functions include \path{asymmetric_odd_derivatives},
\path{feedback_odd_derivatives}, \path{asymmetric_polynomials},
\texttt{section}, \path{phase_lanes}, \path{merge_lanes},
\path{ordinary_taps}, \path{hurwitz_product},
\path{derivative_shift}, \path{integral_shift},
\path{polyphase_tap_matrix}, and \path{apply_polyphase_matrix}.
All core arithmetic is \texttt{fractions.Fraction} or integer
arithmetic. The phase and tap API applies to any finite rational
sequence, not just the special-function catalog.

The final completed run passed \textbf{8,038 exact assertions} and is
recorded in \path{evidence/family_lab_checks.json}. Its scope
includes:

\begin{itemize}
\tightlist
\item
  18 catalog rows through \(A_{12}\), independently compared with dense
  formal Picard iteration through degree 25.
\item
  Bivariate polynomials through \(A_8\), with 540 nonzero monomials
  checked for integer positivity, symmetry, parity and the degree bound.
\item
  Tangent coefficients independently obtained by division of sine and
  cosine, plus hyperbolic tangent, scaled sine and linear anchors.
\item
  Root-filter multipliers reduced exactly in cyclotomic fields for bases
  one through seven and indices zero through 31.
\item
  Direct FIR convolution compared with phase-matrix convolution for 26
  input lengths, seven bases and four tap vectors, including signed
  rational taps.
\item
  The antiderivative interpretation of EGF taps and the factorial
  conversion of Hurwitz products, independently compared with ordinary
  Taylor products.
\item
  Negative controls for omitted phase carries, aliasing lanes by
  concatenation, silent truncation, derivative/delay confusion and
  Hurwitz/FIR confusion.
\end{itemize}

The bounded checks validate the implemented paths in the recorded
ranges. The universal polynomial, degree, clearing, root-filter and
phase-matrix statements follow from the proofs above. The finite checks
do not establish entire continuation, growth order, or historical
novelty. Those claims are separately proved or explicitly left
unresolved in the project.

\path{evidence/family_lab_checks.json} records software version,
code hashes, exact arithmetic, bounds, category counts and run time. It
reports no failed assertion in this run. Independent audit identified
two API input-validation defects: a fractional residue was accepted by
the section operator, and an empty phase matrix accepted base zero. Both
are now rejected and have passing regression checks. These defects
affected invalid inputs, not the displayed identities. A draft prose
example initially transcribed \(A_3(1,2)\) and \(A_4(1,2)\) incorrectly;
the exact catalog comparison exposed the mismatch and corrected them to
\(35200\) and \(38041600\). The code result did not change. Larger
numerical evidence is unnecessary for these algebraic identities; an
independent mathematical audit should examine the reductions.
